\section{Space Complexity --- Window Clock}
%%% generated by the blueprint pipeline from TCSlib.Complexity.SpaceComplexity.WindowClock
\subsection*{Space Complexity --- Window Clock: Introduction}

Declarations of the Lean module \texttt{TCSlib.Complexity.SpaceComplexity.WindowClock}.

\subsection*{Space Complexity --- Window Clock: Results}

\begin{theorem}[Space bounds the work-tape window]
\lean{Turing.FinTM.inWindow_of_spaceUsed_le}
\label{Turing.FinTM.inWindow_of_spaceUsed_le}
\leanok
\uses{Turing.Cfg.InWindow, Turing.FinTM, Turing.MultiTapeTM.abs_pos_lt_card_visited, Turing.MultiTapeTM.initCfg, Turing.MultiTapeTM.mem_visited_of_ne_none, Turing.MultiTapeTM.runFrom, Turing.MultiTapeTM.spaceUsed, Turing.MultiTapeTM.spaceUsedByTape_le_spaceUsed, Turing.MultiTapeTM.visitedByTapeHead, Turing.MultiTapeTM.visitedByTapeHead_mono}
\difficulty{5}
Let $M$ be a finite Turing machine over the Boolean alphabet, let $x$ be a Boolean word,
and let $t,s$ be natural numbers. If the total number of distinct work-tape cells
visited by the computation of $M$ on $x$ through step $t$ is at most $s$, then at every
time $u\le t$ each work-tape head is at a position of absolute value at most $s$, and
every cell at an integer position of absolute value greater than $s$ is blank on every
work tape.
\end{theorem}

\begin{theorem}[Summary of a halting window run]
\lean{Turing.windowRun_summary_of_halts}
\label{Turing.windowRun_summary_of_halts}
\leanok
\uses{Turing.FinTM, Turing.FinTM.ComputesInTime, Turing.FinTM.computesInTime_iff, Turing.FinTM.inWindow_of_spaceUsed_le, Turing.MultiTapeTM.initCfg, Turing.MultiTapeTM.runFrom, Turing.MultiTapeTM.spaceUsed, Turing.outSummary, Turing.outSummaryCode, Turing.summaryTag, Turing.summaryTag_append, Turing.windowRun, Turing.windowRun_halted_of_run}
\difficulty{3}
Let $M$ be a finite-state multi-tape Turing machine with Boolean symbols, and let $x,w$
be Boolean words. Suppose that on input $x$, the machine halts after $t$ steps with
output $w$, using exactly $s_0$ units of space, where $s_0\le s$, and let $t<N$. Then
the windowed run with radius $s$, clock bound $N$, and the output-summary update,
started from the initial configuration on $x$ with the empty summary, reports halting
and returns the configuration reached after $t$ steps together with the code for the
summary of $w$.
\end{theorem}

\begin{theorem}[Configuration bound]
\lean{Turing.FinTM.configBound_le_two_pow}
\label{Turing.FinTM.configBound_le_two_pow}
\leanok
\uses{Turing.FinTM, Turing.FinTM.configBound, Turing.posWidth}
\difficulty{4}
Let $M$ be a finite Turing machine over the Boolean alphabet, with $k$ work tapes and
state set $Q$, and let $n,s\in\mathbb{N}$. If $w(n)$ is the length of the binary
expansion of $n+1$, then $(|Q|+1)(n+2)3^{k(2s+1)}(2s+1)^k\le 2^{|Q|+w(n)+3k(2s+1)}$.
\end{theorem}

\begin{theorem}[A bounded machine for simulating a windowed run]
\lean{Turing.exists_configWindowRunnerTM}
\label{Turing.exists_configWindowRunnerTM}
\leanok
\uses{Turing.Cfg, Turing.Cfg.InWindow, Turing.FinTM, Turing.FinTM.configBound, Turing.MultiTapeTM, Turing.MultiTapeTM.runFrom, Turing.MultiTapeTM.spaceUsed, Turing.RecSym, Turing.RecSym.digitEmb, Turing.RecordAt, Turing.WindowExit, Turing.cfgRecord, Turing.fieldsTape, Turing.recordLength, Turing.windowRun}
Let $M$ be a finite-state Turing machine with Boolean input and $k$ work tapes. Let $T$
be a finite set with decidable equality, and let
$u:T\times\mathrm{Option}(\mathrm{Bool})\to T$ satisfy $u(\tau,\mathrm{none})=\tau$ for
every $\tau\in T$. There exist a natural number $K$, a finite state set $P$ with
decidable equality, a Turing machine $R$ with $K+1$ work tapes over an alphabet
recording Boolean digits, work-tape cells, states, and tags, an entry state
$p_{\mathrm{in}}$, an injective map $e$ from the three window-exit types to states of
$R$ whose values differ from $p_{\mathrm{in}}$, and a natural number $C$ with the
following property. For every Boolean input $x$, natural number $s$, configuration $c$
of $M$ on $x$ whose work-tape head positions have absolute value at most $s$ and whose
work-tape cells outside that window are blank, and tag $\tau\in T$, suppose a
configuration $d$ of $R$ on the digit encoding of $x$ has state $p_{\mathrm{in}}$, has
on its first work tape the record encoding $(c,\tau)$, and has blank contents on every
other work tape. Then there are a time $t$ and an exit type $o$ such that $t\le
C2^{C(\ell+1)}$, where $\ell=k(2s+1)+\mathrm{posWidth}(|x|)+k+4$, and the state of $R$
after $t$ steps is $e(o)$, while before time $t$ it is never in any exit state. At every
time $v\le t$, each work-tape head of $R$ has moved from its initial position by at most
$C(\ell+1)$, and the output remains unchanged; by time $t$ its total space use is at
most $(K+1)(2C(\ell+1)+1)$. For every $N\le \mathrm{configBound}(M,|x|,s)+1$, if the
$N$-step windowed run from $(c,\tau)$ does not end by clock exhaustion, its exit type is
$o$. If $o$ is not clock exhaustion, there is some $N$ for which that windowed run ends
with type $o$, and at time $t$ the runner's first work tape agrees with the record of
the resulting configuration and updated tag throughout the recorded field; its state is
$e(o)$, and all other tape cells retain their initial contents.
\end{theorem}

%%% added by the blueprint pipeline
\subsection*{Space Complexity --- Window Clock: further results}

\begin{lemma}[Value of an all-ones binary word]
\lean{Turing.cw_true_value}
\label{Turing.cw_true_value}
\leanok
\difficulty{4}
For every $W\in\mathbb{N}$, the natural number whose binary digits are $W$ ones equals $2^W-1$.
\end{lemma}

\begin{lemma}[Per-round cost bound]
\lean{Turing.cw_round_cost}
\label{Turing.cw_round_cost}
\leanok
\difficulty{3}
Let $L,n,p\in\mathbb{N}$ with $p\le L$ and $n+2\le 2^p$. Then $L+n+1\le 2^{L+1}$.
\end{lemma}

\begin{lemma}[Clock budget arithmetic]
\lean{Turing.cw_budget}
\label{Turing.cw_budget}
\leanok
\uses{Turing.cw_round_cost}
\difficulty{5}
Let $q,k,n,s,p,L,B,C\in\mathbb{N}$ satisfy $L=k(2s+1)+p+k+4$, $n+2\le 2^p$ and $B\le
2^{q+p+3k(2s+1)}$. Then \[C(B+1)(L+n+1)\le 2^{(C+q+8)(k(2s+1)+p+2)}-1.\]
\end{lemma}

\begin{lemma}[Single-exponential time envelope]
\lean{Turing.cw_time_bound}
\label{Turing.cw_time_bound}
\leanok
\difficulty{5}
Let $c,L,W,A,D\in\mathbb{N}$ with $W\le cL$, and set $C=A+4c+D+10$. Then
\[A(L+1)+\bigl(5(2^W-1)+2W+3\bigr)+(2W+2)+D\le C\cdot 2^{C(L+1)}.\]
\end{lemma}

\begin{lemma}[Linear head envelope with a linear clock]
\lean{Turing.cw_head_bound}
\label{Turing.cw_head_bound}
\leanok
\difficulty{3}
Let $c,L,W,C\in\mathbb{N}$ with $W\le cL$. Then $C(L+W+1)\le C(c+1)(L+1)$.
\end{lemma}

\begin{lemma}[Total field length of a configuration record]
\lean{Turing.cw_field_total}
\label{Turing.cw_field_total}
\leanok
\uses{Turing.Cfg, Turing.cfgRecord, Turing.cfgRecord_length, Turing.joinFields_cfgRecord_length, Turing.joinFields_length, Turing.posWidth, Turing.recordLength_eq}
\difficulty{4}
Let $k\in\mathbb{N}$, let $x$ be a Boolean input, let $s\in\mathbb{N}$, let $d$ be a
configuration with $k$ work tapes on input $x$, and let $\tau$ be a tag. Consider the
configuration record of $d$ at radius $s$ with tag $\tau$: the $k+3$ fields consisting of the
state, the input position in binary with as many digits as the bit length $w$ of $|x|+1$, the
cells $-s,\dots,s$ of each work tape with head marks, and the tag. The sum of the lengths of
these fields is $k(2s+1)+w+2$.
\end{lemma}

\begin{lemma}[An all-ones clock covers every round]
\lean{Turing.cw_record_clock}
\label{Turing.cw_record_clock}
\leanok
\uses{Turing.Cfg, Turing.FinTM, Turing.FinTM.configBound, Turing.FinTM.configBound_le_two_pow, Turing.cfgRecord, Turing.cw_budget, Turing.cw_field_total, Turing.cw_true_value, Turing.posWidth, Turing.recordLength, Turing.recordLength_eq, Turing.record_pos_lt}
\difficulty{4}
Let $M$ be a finite Turing machine over the Boolean alphabet with $k$ work tapes and state set
$Q$, let $x$ be a Boolean input of length $n$, let $s,C\in\mathbb{N}$, let $d$ be a
configuration of $M$ on $x$ and $\tau$ a tag. Let $w$ be the bit length of $n+1$, let
$\ell=k(2s+1)+w+k+4$ be the record length, let $F$ be the total length of the fields of the
configuration record of $d$ at radius $s$ with tag $\tau$, and set $c=C+|Q|+8$ and $W=cF$. Then
$W\le c\ell$, and for every $N\le (|Q|+1)(n+2)3^{k(2s+1)}(2s+1)^k+1$ the value $2^W-1$ of the
binary word of $W$ ones satisfies $CN(\ell+n+1)\le 2^W-1$.
\end{lemma}

\begin{definition}[One-step left move on a tape]
\lean{Turing.cw_leftTM}
\label{Turing.cw_leftTM}
\leanok
\uses{Turing.MultiTapeTM, Turing.tapeAction}
For $k\in\mathbb{N}$ and a work tape $i<k$, the one-step left move on tape $i$ is the Turing
machine with $k$ work tapes and Boolean states, starting in state false, whose every transition
keeps the input head, writes nothing, moves the head of tape $i$ one cell to the left, leaves
every other tape and the output untouched, and enters state true.
\end{definition}

\begin{lemma}[Run of the one-step left move]
\lean{Turing.cw_left_run}
\label{Turing.cw_left_run}
\leanok
\uses{Turing.Cfg, Turing.MultiTapeTM.runFrom, Turing.MultiTapeTM.runFrom_one, Turing.MultiTapeTM.runFrom_zero, Turing.cw_leftTM, Turing.moveInputPos_zero, Turing.step_of_tapeAction_move}
\difficulty{3}
Let $k\in\mathbb{N}$, let $i<k$, and consider the Turing machine with $k$ work tapes and
Boolean states whose every transition moves the head of tape $i$ one cell left, writes and
emits nothing, keeps the input head, and enters state true. Let $d$ be a configuration of this
machine whose state is false. Then after one step the configuration is $d$ with state true and
the head of tape $i$ moved from position $p$ to $p-1$, and at time $0$ the state is not true.
\end{lemma}

\begin{lemma}[Reading a constant patch]
\lean{Turing.cw_patch_rep}
\label{Turing.cw_patch_rep}
\leanok
\uses{Turing.FinTM.bufferTape_nat, Turing.wordPatch, Turing.wordPatch_of_mem, Turing.wordPatch_of_not_mem}
\difficulty{4}
Let $\Gamma$ be an alphabet, let $\gamma$ be a tape over $\Gamma$ (a map from $\mathbb{Z}$ to
cells that are blank or hold a letter), and let $p\in\mathbb{Z}$, $n\in\mathbb{N}$ and
$e\in\Gamma$. The tape obtained from $\gamma$ by writing $n$ copies of $e$ in the cells
$p,\dots,p+n-1$ holds $e$ at every cell $z$ with $p\le z<p+n$, and holds $\gamma(z)$ at every
other cell $z$.
\end{lemma}

\begin{lemma}[Tape holding a constant word]
\lean{Turing.cw_buffer_rep}
\label{Turing.cw_buffer_rep}
\leanok
\uses{Turing.FinTM.bufferTape, Turing.FinTM.bufferTape_nat, Turing.FinTM.bufferTape_of_neg}
\difficulty{4}
Let $n\in\mathbb{N}$, let $e$ be a letter and let $z\in\mathbb{Z}$. The tape that holds the
word $e^n$ in cells $0,\dots,n-1$ and is blank elsewhere holds $e$ at cell $z$ if $0\le z<n$,
and is blank at $z$ otherwise.
\end{lemma}

\begin{lemma}[Joining adjacent constant patches]
\lean{Turing.cw_join}
\label{Turing.cw_join}
\leanok
\uses{Turing.wordPatch, Turing.wordPatch_append}
\difficulty{2}
Let $\gamma$ be a tape over an alphabet $\Gamma$, let $p\in\mathbb{Z}$, $e\in\Gamma$ and
$n,l\in\mathbb{N}$. Writing $e^n$ onto $\gamma$ from cell $p$ and then $e^l$ from cell $p+n$
gives the same tape as writing $e^{n+l}$ onto $\gamma$ from cell $p$.
\end{lemma}

\begin{lemma}[Constant mapped copy onto a destination tape]
\lean{Turing.cw_copy_constant}
\label{Turing.cw_copy_constant}
\leanok
\uses{Turing.Cfg, Turing.FinTM.bufferTape, Turing.FinTM.bufferTape_nat, Turing.MultiTapeTM.runFrom, Turing.SweepPhase, Turing.copyMapTM, Turing.copyMapTM_run_ofCfg, Turing.cw_patch_rep, Turing.wordPatch}
\difficulty{5}
Let $k\in\mathbb{N}$, let $\sigma\ne\delta$ be two work tapes among $k$, let $e$ be a letter
and $w$ a word. Consider the mapped copy machine from tape $\sigma$ to tape $\delta$ that
writes $e$ in place of every copied letter, and let $d$ be a configuration of it in its sweep
state such that for every integer $-1\le z\le|w|$ the cell of tape $\sigma$ at offset $z$ from
its head holds the $z$-th letter of $w$ (blank at $z=-1$ and $z=|w|$). Then after $2|w|+2$
steps the machine is in its done state and its configuration is $d$ except that $e^{|w|}$ has
been written on tape $\delta$ in the $|w|$ cells starting at its head; it is not done at any
earlier time; and at every time up to $2|w|+2$ the heads of $\sigma$ and $\delta$ each lie in
$[h-1,h+|w|]$, $h$ being that head's initial position, while every other head is at its initial
position.
\end{lemma}

\begin{lemma}[Clearing a written patch blanks the tape]
\lean{Turing.cw_clear_patch}
\label{Turing.cw_clear_patch}
\leanok
\uses{Turing.Cfg, Turing.FinTM.bufferTape, Turing.FinTM.bufferTape_left, Turing.FinTM.bufferTape_nat, Turing.MultiTapeTM.runFrom, Turing.SweepPhase, Turing.clearTM, Turing.clearTM_run_ofCfg, Turing.wordPatch, Turing.wordPatch_of_mem, Turing.wordPatch_of_not_mem}
\difficulty{5}
Let $k\in\mathbb{N}$, let $i$ be one of $k$ work tapes, let $w$ be a word, and let $d$ be a
configuration whose tape $i$ is blank at every cell outside $[h,h+|w|)$, where $h$ is the head
position of tape $i$ in $d$. Let $A$ be $d$ with state sweep and with $w$ written on tape $i$
from cell $h$. Then the clear machine on tape $i$, run from $A$ for $2|w|+2$ steps, reaches $d$
with state done and tape $i$ entirely blank; it is not in state done at any earlier time; and
at every time up to $2|w|+2$ the head of tape $i$ lies in $[h-1,h+|w|]$ while every other head
stays at its position in $d$.
\end{lemma}

\begin{definition}[Phases of the clock runner]
\lean{Turing.cw_Idx}
\label{Turing.cw_Idx}
\leanok
For $c,m\in\mathbb{N}$, the phases of the clock runner form the following finite set with
decidable equality: for each pass $p<c$ and field $j<m$, the copy of field $j$ onto the clock,
the skip past field $j$ on the record tape, the one-cell move back onto the clock's delimiter,
and the skip back over field $j$ on the record tape; the skip past the written clock cells,
indexed either by a pair $(p,j)$ or by a marker for the end of the last pass; the skip back
over the clock word; the windowed runner; the clear of the clock, tagged by a Boolean recording
a halted or a left exit; and the stop phase.
\end{definition}

\begin{definition}[Control sets of the clock runner's phases]
\lean{Turing.cw_Ctrl}
\label{Turing.cw_Ctrl}
\leanok
\uses{Turing.SeekPhase, Turing.SweepPhase, Turing.WindowExit, Turing.cw_Idx}
Given a state set $P$ and $c,m\in\mathbb{N}$, each phase of the clock runner is assigned a
control set: the sweep controls (sweep, rewind, done) for the field copies and the clears; the
seek controls (run, gap, back, done) for the record skips, the clock skips, the record
skip-backs and the clock skip-back; the Booleans for the one-cell left moves; $P$ for the
windowed runner; and the three window exits (halted, left, exhausted) for the stop phase.
\end{definition}

\begin{definition}[Machines of the clock runner's phases]
\lean{Turing.cw_rows}
\label{Turing.cw_rows}
\leanok
\uses{Turing.MultiTapeTM, Turing.clearTM, Turing.copyMapTM, Turing.cw_Ctrl, Turing.cw_Idx, Turing.cw_leftTM, Turing.skipBackTM, Turing.skipTM}
Let $K\in\mathbb{N}$, let $e$ be a tape symbol, let $R$ be a Turing machine with $K+2$ work
tapes and state set $P$, and let $c,m\in\mathbb{N}$. Each phase of the clock runner is assigned
a machine on $K+2$ work tapes: each field copy runs the mapped copy from tape $0$ to tape $K+1$
that writes $e$ in every cell; each record skip runs the skip on tape $0$; each clock skip runs
the skip on tape $K+1$; each one-cell left move moves the head of tape $K+1$ one cell left;
each record skip-back runs the skip back on tape $0$; the clock skip-back runs the skip back on
tape $K+1$; the runner phase runs $R$; each clear runs the clear of tape $K+1$; and the stop
phase is the stationary machine, started at the exit exhausted, that keeps its state and moves,
writes and emits nothing.
\end{definition}

\begin{definition}[Continuation after a field's round]
\lean{Turing.cw_next}
\label{Turing.cw_next}
\leanok
\uses{Turing.SeekPhase, Turing.SweepPhase, Turing.cw_Ctrl, Turing.cw_Idx}
Let $c,m\in\mathbb{N}$, $p<c$ and $j<m$. The continuation of the clock runner after the round
of field $j$ in pass $p$ is the copy phase of field $j+1$ in pass $p$ with control sweep when
$j+1<m$, and otherwise the skip back over field $j$ in pass $p$ with control run.
\end{definition}

\begin{definition}[Continuation after a record skip-back]
\lean{Turing.cw_prev}
\label{Turing.cw_prev}
\leanok
\uses{Turing.SeekPhase, Turing.SweepPhase, Turing.cw_Ctrl, Turing.cw_Idx}
Let $c,m\in\mathbb{N}$, $p<c$ and $j<m$. The continuation of the clock runner after the skip
back over field $j$ in pass $p$ is the skip back over field $j-1$ in pass $p$ with control run
when $j>0$; otherwise the copy phase of field $0$ in pass $p+1$ with control sweep when
$p+1<c$; and otherwise the final clock skip (after the last pass) with control run.
\end{definition}

\begin{definition}[Continuation after a clock skip]
\lean{Turing.cw_afterSkip}
\label{Turing.cw_afterSkip}
\leanok
\uses{Turing.SeekPhase, Turing.cw_Ctrl, Turing.cw_Idx}
Let $c,m\in\mathbb{N}$. After the clock skip indexed by a pair $(p,j)$ with $p<c$ and $j<m$,
the clock runner continues with the one-cell left move of $(p,j)$ in control false; after the
final clock skip, which carries no pair, it continues with the skip back over the clock word in
control run.
\end{definition}

\begin{definition}[Redirect table of the clock runner]
\lean{Turing.cw_red}
\label{Turing.cw_red}
\leanok
\uses{Turing.SeekPhase, Turing.SweepPhase, Turing.WindowExit, Turing.cw_Ctrl, Turing.cw_Idx, Turing.cw_afterSkip, Turing.cw_next, Turing.cw_prev}
Let $P$ be a state set with decidable equality, $p_{\mathrm{en}}\in P$ an entry state,
$\mathrm{ex}$ a map from the window exits (halted, left, exhausted) to $P$, and
$c,m\in\mathbb{N}$. The redirect table sends the copy of field $j$ in pass $p$ at done to the
record skip of $(p,j)$ at run; that record skip at done to the clock skip indexed by $(p,j)$ at
run; a clock skip at done to its continuation (the left move of $(p,j)$ at false, or after the
last pass the clock skip-back at run); the left move of $(p,j)$ at true to the copy of field
$j+1$ in pass $p$ at sweep if $j+1<m$, else to the skip back over field $j$ at run; the skip
back over field $j$ in pass $p$ at done to the skip back over field $j-1$ at run if $j>0$, else
to the copy of field $0$ in pass $p+1$ at sweep if $p+1<c$, else to the final clock skip at
run; the clock skip-back at done to the runner at $p_{\mathrm{en}}$; the runner at
$\mathrm{ex}(\text{halted})$, $\mathrm{ex}(\text{left})$ and $\mathrm{ex}(\text{exhausted})$ to
the clear tagged true at sweep, the clear tagged false at sweep, and the stop phase at
exhausted, respectively; and the clear tagged $b$ at done to the stop phase at halted if $b$ is
true and at left otherwise. Every other phase and control is not redirected.
\end{definition}

\begin{definition}[The clock runner machine]
\lean{Turing.cw_hostTM}
\label{Turing.cw_hostTM}
\leanok
\uses{Turing.MultiTapeTM, Turing.WindowExit, Turing.cw_Ctrl, Turing.cw_Idx, Turing.cw_red, Turing.cw_rows, Turing.phaseHostTM}
Let $K,c,m\in\mathbb{N}$, let $e$ be a tape symbol, and let $R$ be a Turing machine with $K+2$
work tapes whose state set $P$ has decidable equality, with entry state $p_{\mathrm{en}}\in P$
and a map $\mathrm{ex}$ from the window exits to $P$. The clock runner is the Turing machine
with $K+2$ work tapes whose states are pairs of a phase and a control of that phase: at a pair
that the clock runner's redirect table sends to another pair, it takes one stationary, silent
step to that pair; at any other pair it takes the transition of the phase's machine (field
copies, record and clock skips, left moves, skip-backs, the runner $R$, the clear of tape
$K+1$, or the stationary stop machine), tagging the new control with the same phase. Its
initial state is the stop phase at exhausted.
\end{definition}

\begin{definition}[Writer configurations]
\lean{Turing.cw_cfg}
\label{Turing.cw_cfg}
\leanok
\uses{Turing.Cfg, Turing.wordPatch}
Let $K\in\mathbb{N}$, let $e$ be a tape symbol, let $d$ be a configuration with $K+2$ work
tapes on an input $x$, let $q$ be a state, let $a,b\in\mathbb{Z}$ and $n\in\mathbb{N}$. The
writer configuration has state $q$, the input head of $d$, the output of $d$, tape $K+1$ equal
to that of $d$ with $n$ copies of $e$ written from the head position of tape $K+1$ in $d$,
every other tape as in $d$, the head of tape $0$ moved by $a$ and the head of tape $K+1$ moved
by $b$ from their positions in $d$, and every other head as in $d$.
\end{definition}

\begin{lemma}[Record and clock tapes are distinct]
\lean{Turing.cw_zero_ne_last}
\label{Turing.cw_zero_ne_last}
\leanok
\difficulty{3}
For every $K\in\mathbb{N}$, the indices $0$ and $K+1$ of the work tapes $\{0,1,\dots,K+1\}$ are
distinct.
\end{lemma}

\begin{lemma}[Clock tape of a writer configuration]
\lean{Turing.cw_cfg_clockTape}
\label{Turing.cw_cfg_clockTape}
\leanok
\uses{Turing.Cfg, Turing.cw_cfg, Turing.wordPatch}
\difficulty{1}
Let $K\in\mathbb{N}$, let $d$ be a configuration of a Turing machine with $K+2$ work tapes on
some input, and let $e$ be a tape symbol. For a state $q$, integers $a,b$ and $n\in\mathbb{N}$,
let $\omega(q,a,n,b)$ be the configuration obtained from $d$ by setting the state to $q$,
moving the head of tape $0$ by $a$ and the head of tape $K+1$ by $b$, and writing $n$ copies of
$e$ on tape $K+1$ starting at its head position in $d$; let $h_j$ denote the head position of
tape $j$ in $d$. Then tape $K+1$ of $\omega(q,a,n,b)$ is tape $K+1$ of $d$ with $n$ copies of
$e$ written in the cells starting at $h_{K+1}$.
\end{lemma}

\begin{lemma}[Clock cell of a writer configuration]
\lean{Turing.cw_cfg_clock}
\label{Turing.cw_cfg_clock}
\leanok
\uses{Turing.Cfg, Turing.cw_cfg, Turing.cw_cfg_clockTape, Turing.cw_patch_rep}
\difficulty{1}
Let $K\in\mathbb{N}$, let $d$ be a configuration of a Turing machine with $K+2$ work tapes on
some input, and let $e$ be a tape symbol. For a state $q$, integers $a,b$ and $n\in\mathbb{N}$,
let $\omega(q,a,n,b)$ be the configuration obtained from $d$ by setting the state to $q$,
moving the head of tape $0$ by $a$ and the head of tape $K+1$ by $b$, and writing $n$ copies of
$e$ on tape $K+1$ starting at its head position in $d$; let $h_j$ denote the head position of
tape $j$ in $d$. Then for every $z\in\mathbb{Z}$, cell $z$ of tape $K+1$ in $\omega(q,a,n,b)$
holds $e$ if $h_{K+1}\le z<h_{K+1}+n$, and otherwise holds the content of cell $z$ of tape
$K+1$ in $d$.
\end{lemma}

\begin{lemma}[Other tapes of a writer configuration]
\lean{Turing.cw_cfg_tape}
\label{Turing.cw_cfg_tape}
\leanok
\uses{Turing.Cfg, Turing.cw_cfg}
\difficulty{1}
Let $K\in\mathbb{N}$, let $d$ be a configuration of a Turing machine with $K+2$ work tapes on
some input, and let $e$ be a tape symbol. For a state $q$, integers $a,b$ and $n\in\mathbb{N}$,
let $\omega(q,a,n,b)$ be the configuration obtained from $d$ by setting the state to $q$,
moving the head of tape $0$ by $a$ and the head of tape $K+1$ by $b$, and writing $n$ copies of
$e$ on tape $K+1$ starting at its head position in $d$; let $h_j$ denote the head position of
tape $j$ in $d$. Then for every work tape $j\ne K+1$ and every cell $z$, cell $z$ of tape $j$
in $\omega(q,a,n,b)$ equals cell $z$ of tape $j$ in $d$.
\end{lemma}

\begin{lemma}[Record head of a writer configuration]
\lean{Turing.cw_cfg_pos₀}
\label{Turing.cw_cfg_pos₀}
\leanok
\uses{Turing.Cfg, Turing.cw_cfg}
\difficulty{1}
Let $K\in\mathbb{N}$, let $d$ be a configuration of a Turing machine with $K+2$ work tapes on
some input, and let $e$ be a tape symbol. For a state $q$, integers $a,b$ and $n\in\mathbb{N}$,
let $\omega(q,a,n,b)$ be the configuration obtained from $d$ by setting the state to $q$,
moving the head of tape $0$ by $a$ and the head of tape $K+1$ by $b$, and writing $n$ copies of
$e$ on tape $K+1$ starting at its head position in $d$; let $h_j$ denote the head position of
tape $j$ in $d$. Then the head of tape $0$ in $\omega(q,a,n,b)$ is at position $h_0+a$.
\end{lemma}

\begin{lemma}[Clock head of a writer configuration]
\lean{Turing.cw_cfg_posL}
\label{Turing.cw_cfg_posL}
\leanok
\uses{Turing.Cfg, Turing.cw_cfg, Turing.cw_zero_ne_last}
\difficulty{1}
Let $K\in\mathbb{N}$, let $d$ be a configuration of a Turing machine with $K+2$ work tapes on
some input, and let $e$ be a tape symbol. For a state $q$, integers $a,b$ and $n\in\mathbb{N}$,
let $\omega(q,a,n,b)$ be the configuration obtained from $d$ by setting the state to $q$,
moving the head of tape $0$ by $a$ and the head of tape $K+1$ by $b$, and writing $n$ copies of
$e$ on tape $K+1$ starting at its head position in $d$; let $h_j$ denote the head position of
tape $j$ in $d$. Then the head of tape $K+1$ in $\omega(q,a,n,b)$ is at position $h_{K+1}+b$.
\end{lemma}

\begin{lemma}[Other heads of a writer configuration]
\lean{Turing.cw_cfg_pos}
\label{Turing.cw_cfg_pos}
\leanok
\uses{Turing.Cfg, Turing.cw_cfg}
\difficulty{1}
Let $K\in\mathbb{N}$, let $d$ be a configuration of a Turing machine with $K+2$ work tapes on
some input, and let $e$ be a tape symbol. For a state $q$, integers $a,b$ and $n\in\mathbb{N}$,
let $\omega(q,a,n,b)$ be the configuration obtained from $d$ by setting the state to $q$,
moving the head of tape $0$ by $a$ and the head of tape $K+1$ by $b$, and writing $n$ copies of
$e$ on tape $K+1$ starting at its head position in $d$; let $h_j$ denote the head position of
tape $j$ in $d$. Then for every work tape $j$ with $j\ne 0$ and $j\ne K+1$, the head of tape
$j$ in $\omega(q,a,n,b)$ is at position $h_j$.
\end{lemma}

\begin{lemma}[State renaming keeps the heads]
\lean{Turing.cw_mapState_pos}
\label{Turing.cw_mapState_pos}
\leanok
\uses{Turing.Cfg, Turing.Cfg.mapState}
\difficulty{1}
Let $C$ be a configuration of a multi-tape Turing machine with state set $S$, and let $f:S\to
S'$. The configuration obtained from $C$ by renaming its state through $f$ has the same
work-tape head positions as $C$.
\end{lemma}

\begin{lemma}[State renaming keeps the tapes]
\lean{Turing.cw_mapState_tapes}
\label{Turing.cw_mapState_tapes}
\leanok
\uses{Turing.Cfg, Turing.Cfg.mapState}
\difficulty{1}
Let $C$ be a configuration of a multi-tape Turing machine with state set $S$, and let $f:S\to
S'$. The configuration obtained from $C$ by renaming its state through $f$ has the same
work-tape contents as $C$.
\end{lemma}

\begin{lemma}[State renaming keeps input head and output]
\lean{Turing.cw_mapState_io}
\label{Turing.cw_mapState_io}
\leanok
\uses{Turing.Cfg, Turing.Cfg.mapState}
\difficulty{1}
Let $C$ be a configuration of a multi-tape Turing machine with state set $S$, and let $f:S\to
S'$. The configuration obtained from $C$ by renaming its state through $f$ has the same
input-head position and the same output as $C$.
\end{lemma}

\begin{lemma}[Renaming the state of a writer configuration]
\lean{Turing.cw_cfg_mapState}
\label{Turing.cw_cfg_mapState}
\leanok
\uses{Turing.Cfg, Turing.Cfg.mapState, Turing.cw_cfg, Turing.cw_mapState_io, Turing.cw_mapState_pos, Turing.cw_mapState_tapes}
\difficulty{3}
Let $K\in\mathbb{N}$, let $d$ be a configuration of a Turing machine with $K+2$ work tapes on
some input, and let $e$ be a tape symbol. For a state $q$, integers $a,b$ and $n\in\mathbb{N}$,
let $\omega(q,a,n,b)$ be the configuration obtained from $d$ by setting the state to $q$,
moving the head of tape $0$ by $a$ and the head of tape $K+1$ by $b$, and writing $n$ copies of
$e$ on tape $K+1$ starting at its head position in $d$; let $h_j$ denote the head position of
tape $j$ in $d$. Then for every map $f$ on states, renaming the state of $\omega(q,a,n,b)$
through $f$ gives $\omega(f(q),a,n,b)$.
\end{lemma}

\begin{lemma}[Unmoved writer configuration keeps the heads]
\lean{Turing.cw_cfg_pos_zero}
\label{Turing.cw_cfg_pos_zero}
\leanok
\uses{Turing.Cfg, Turing.cw_cfg, Turing.cw_cfg_pos, Turing.cw_cfg_posL, Turing.cw_cfg_pos₀}
\difficulty{4}
Let $K\in\mathbb{N}$, let $d$ be a configuration of a Turing machine with $K+2$ work tapes on
some input, and let $e$ be a tape symbol. For a state $q$, integers $a,b$ and $n\in\mathbb{N}$,
let $\omega(q,a,n,b)$ be the configuration obtained from $d$ by setting the state to $q$,
moving the head of tape $0$ by $a$ and the head of tape $K+1$ by $b$, and writing $n$ copies of
$e$ on tape $K+1$ starting at its head position in $d$; let $h_j$ denote the head position of
tape $j$ in $d$. Then for every state $q$ and every $n\in\mathbb{N}$, the work-tape head
positions of $\omega(q,0,n,0)$ are those of $d$.
\end{lemma}

\begin{lemma}[Trivial writer configuration is the start]
\lean{Turing.cw_cfg_zero}
\label{Turing.cw_cfg_zero}
\leanok
\uses{Turing.Cfg, Turing.cw_cfg, Turing.cw_cfg_clock, Turing.cw_cfg_pos_zero, Turing.cw_cfg_tape}
\difficulty{4}
Let $K\in\mathbb{N}$, let $d$ be a configuration of a Turing machine with $K+2$ work tapes on
some input, and let $e$ be a tape symbol. For a state $q$, integers $a,b$ and $n\in\mathbb{N}$,
let $\omega(q,a,n,b)$ be the configuration obtained from $d$ by setting the state to $q$,
moving the head of tape $0$ by $a$ and the head of tape $K+1$ by $b$, and writing $n$ copies of
$e$ on tape $K+1$ starting at its head position in $d$; let $h_j$ denote the head position of
tape $j$ in $d$. If the state of $d$ is $q$, then $\omega(q,0,0,0)=d$.
\end{lemma}

\begin{lemma}[Moving the record head of a writer configuration]
\lean{Turing.cw_cfg_move₀}
\label{Turing.cw_cfg_move₀}
\leanok
\uses{Turing.Cfg, Turing.cw_cfg, Turing.cw_cfg_pos₀}
\difficulty{4}
Let $K\in\mathbb{N}$, let $d$ be a configuration of a Turing machine with $K+2$ work tapes on
some input, and let $e$ be a tape symbol. For a state $q$, integers $a,b$ and $n\in\mathbb{N}$,
let $\omega(q,a,n,b)$ be the configuration obtained from $d$ by setting the state to $q$,
moving the head of tape $0$ by $a$ and the head of tape $K+1$ by $b$, and writing $n$ copies of
$e$ on tape $K+1$ starting at its head position in $d$; let $h_j$ denote the head position of
tape $j$ in $d$. Let $q'$ be a state and $v,a'\in\mathbb{Z}$ with $v=h_0+a'$. Then the
configuration obtained from $\omega(q,a,n,b)$ by setting the state to $q'$ and placing the head
of tape $0$ at $v$ is $\omega(q',a',n,b)$.
\end{lemma}

\begin{lemma}[Moving the clock head of a writer configuration]
\lean{Turing.cw_cfg_moveL}
\label{Turing.cw_cfg_moveL}
\leanok
\uses{Turing.Cfg, Turing.cw_cfg, Turing.cw_cfg_pos, Turing.cw_cfg_posL, Turing.cw_cfg_pos₀}
\difficulty{4}
Let $K\in\mathbb{N}$, let $d$ be a configuration of a Turing machine with $K+2$ work tapes on
some input, and let $e$ be a tape symbol. For a state $q$, integers $a,b$ and $n\in\mathbb{N}$,
let $\omega(q,a,n,b)$ be the configuration obtained from $d$ by setting the state to $q$,
moving the head of tape $0$ by $a$ and the head of tape $K+1$ by $b$, and writing $n$ copies of
$e$ on tape $K+1$ starting at its head position in $d$; let $h_j$ denote the head position of
tape $j$ in $d$. Let $q'$ be a state and $v,b'\in\mathbb{Z}$ with $v=h_{K+1}+b'$. Then the
configuration obtained from $\omega(q,a,n,b)$ by setting the state to $q'$ and placing the head
of tape $K+1$ at $v$ is $\omega(q',a,n,b')$.
\end{lemma}

\begin{definition}[Head box around an origin]
\lean{Turing.cw_box}
\label{Turing.cw_box}
\leanok
For $k,r\in\mathbb{N}$ and an origin $o\in\mathbb{Z}^k$, a vector of head positions
$\mathrm{pos}\in\mathbb{Z}^k$ lies in the box of radius $r$ around $o$ when
$\abs{\mathrm{pos}_j-o_j}\le r$ for every $j<k$.
\end{definition}

\begin{lemma}[Box from record and clock head ranges]
\lean{Turing.cw_box_of}
\label{Turing.cw_box_of}
\leanok
\uses{Turing.cw_box}
\difficulty{4}
Let $K,r\in\mathbb{N}$, let $o,\mathrm{pos}\in\mathbb{Z}^{K+2}$, and let
$\ell_0,u_0,\ell_1,u_1\in\mathbb{Z}$. Suppose $o_0+\ell_0\le \mathrm{pos}_0\le o_0+u_0$,
$o_{K+1}+\ell_1\le\mathrm{pos}_{K+1}\le o_{K+1}+u_1$, $\mathrm{pos}_j=o_j$ for every
$j\notin\{0,K+1\}$, and $-r\le\ell_0$, $u_0\le r$, $-r\le\ell_1$, $u_1\le r$. Then
$\abs{\mathrm{pos}_j-o_j}\le r$ for every $j$.
\end{lemma}

\begin{lemma}[Box from one moving head]
\lean{Turing.cw_box_row}
\label{Turing.cw_box_row}
\leanok
\uses{Turing.cw_box}
\difficulty{4}
Let $k,r\in\mathbb{N}$, let $o,s,\mathrm{pos},\ell,u\in\mathbb{Z}^k$, and let $i<k$. Suppose
$\ell_i\le\mathrm{pos}_i\le u_i$, $\mathrm{pos}_j=s_j$ for every $j\ne i$, $o_i-r\le\ell_i$ and
$u_i\le o_i+r$, and $\abs{s_j-o_j}\le r$ for every $j$. Then $\abs{\mathrm{pos}_j-o_j}\le r$
for every $j$.
\end{lemma}

\begin{lemma}[Writer configurations stay in the box]
\lean{Turing.cw_cfg_box}
\label{Turing.cw_cfg_box}
\leanok
\uses{Turing.Cfg, Turing.cw_box, Turing.cw_cfg, Turing.cw_cfg_pos, Turing.cw_cfg_posL, Turing.cw_cfg_pos₀}
\difficulty{4}
Let $K\in\mathbb{N}$, let $d$ be a configuration of a Turing machine with $K+2$ work tapes on
some input, and let $e$ be a tape symbol. For a state $q$, integers $a,b$ and $n\in\mathbb{N}$,
let $\omega(q,a,n,b)$ be the configuration obtained from $d$ by setting the state to $q$,
moving the head of tape $0$ by $a$ and the head of tape $K+1$ by $b$, and writing $n$ copies of
$e$ on tape $K+1$ starting at its head position in $d$; let $h_j$ denote the head position of
tape $j$ in $d$. Let $r\in\mathbb{N}$. If $\abs{a}\le r$ and $\abs{b}\le r$, then every
work-tape head of $\omega(q,a,n,b)$ is within distance $r$ of its position in $d$.
\end{lemma}

\begin{lemma}[No redirect before a phase's anchor]
\lean{Turing.cw_guard}
\label{Turing.cw_guard}
\leanok
\uses{Turing.Cfg, Turing.MultiTapeTM, Turing.MultiTapeTM.runFrom, Turing.WindowExit, Turing.cw_Ctrl, Turing.cw_Idx, Turing.cw_red, Turing.cw_rows}
\difficulty{1}
Let $K,c,m\in\mathbb{N}$, let $e$ be a tape symbol, and let $R$ be a Turing machine with $K+2$
work tapes whose state set $P$ has decidable equality, together with an entry state
$p_{\mathrm{en}}\in P$ and a map $\mathrm{ex}$ from the three window exits (halted, left,
exhausted) to $P$. Let $H$ be the clock runner built from these data: the phase host on $K+2$
work tapes whose phases are the clock writer's phases for $c$ passes over $m$ record fields
(acting on the record tape $0$ and the clock tape $K+1$), the runner $R$, the clear of the
clock tape, and a stationary stop phase; the states of the stop phase are called exits. For a
phase $i$, let $M_i$ be its machine and $\rho_i$ its row of the redirect table. Let $a$ be a
configuration of $M_i$, $t\in\mathbb{N}$ and $v$ a control of phase $i$. Suppose the run of
$M_i$ from $a$ is not in state $v$ at any time before $t$, and $\rho_i$ redirects no control
other than $v$. Then at every time $u<t$, the control of that run is not redirected by
$\rho_i$.
\end{lemma}

\begin{lemma}[One phase as a host segment]
\lean{Turing.cw_seg}
\label{Turing.cw_seg}
\leanok
\uses{Turing.Cfg, Turing.Cfg.mapState, Turing.MultiTapeTM, Turing.MultiTapeTM.runFrom, Turing.Segment, Turing.WindowExit, Turing.cw_Ctrl, Turing.cw_Idx, Turing.cw_hostTM, Turing.cw_mapState_pos, Turing.cw_red, Turing.cw_rows, Turing.phaseHostTM_run, Turing.segment_iff}
\difficulty{5}
Let $K,c,m\in\mathbb{N}$, let $e$ be a tape symbol, and let $R$ be a Turing machine with $K+2$
work tapes whose state set $P$ has decidable equality, together with an entry state
$p_{\mathrm{en}}\in P$ and a map $\mathrm{ex}$ from the three window exits (halted, left,
exhausted) to $P$. Let $H$ be the clock runner built from these data: the phase host on $K+2$
work tapes whose phases are the clock writer's phases for $c$ passes over $m$ record fields
(acting on the record tape $0$ and the clock tape $K+1$), the runner $R$, the clear of the
clock tape, and a stationary stop phase; the states of the stop phase are called exits. For a
phase $i$, let $M_i$ be its machine and $\rho_i$ its row of the redirect table. Let $i$ be a
phase other than the stop phase, let $a,b$ be configurations of $M_i$ such that $M_i$ takes $a$
to $b$ in $t$ steps, let $b$ have control $v$ with $\rho_i(v)=q$, suppose no control of the run
at a time before $t$ is redirected by $\rho_i$, and let $B$ be a property of head vectors that
holds at every time up to $t$ of the run. Then the run of $H$ from $a$ with its state tagged by
phase $i$ reaches $b$ with its state replaced by $q$ after exactly $t+1$ steps, is in no exit
state before time $t+1$, and its head vector satisfies $B$ at every time up to $t+1$.
\end{lemma}

\begin{lemma}[A one-head row as a host segment]
\lean{Turing.cw_seg_row}
\label{Turing.cw_seg_row}
\leanok
\uses{Turing.Cfg, Turing.MultiTapeTM, Turing.MultiTapeTM.runFrom, Turing.Segment, Turing.WindowExit, Turing.cw_Ctrl, Turing.cw_Idx, Turing.cw_box, Turing.cw_box_row, Turing.cw_cfg, Turing.cw_cfg_box, Turing.cw_cfg_mapState, Turing.cw_guard, Turing.cw_hostTM, Turing.cw_red, Turing.cw_rows, Turing.cw_seg}
\difficulty{3}
Let $K,c,m\in\mathbb{N}$, let $e$ be a tape symbol, and let $R$ be a Turing machine with $K+2$
work tapes whose state set $P$ has decidable equality, together with an entry state
$p_{\mathrm{en}}\in P$ and a map $\mathrm{ex}$ from the three window exits (halted, left,
exhausted) to $P$. Let $H$ be the clock runner built from these data: the phase host on $K+2$
work tapes whose phases are the clock writer's phases for $c$ passes over $m$ record fields
(acting on the record tape $0$ and the clock tape $K+1$), the runner $R$, the clear of the
clock tape, and a stationary stop phase; the states of the stop phase are called exits. For a
configuration $d$ of $H$, a control $q$, integers $a,b$ and $n\in\mathbb{N}$, let
$\omega(q,a,n,b)$ be the configuration obtained from $d$ by setting the state to $q$, moving
the head of tape $0$ by $a$ and the head of tape $K+1$ by $b$, and writing $n$ copies of $e$ on
tape $K+1$ starting at its head position in $d$; let $h_j$ denote the head position of tape $j$
in $d$. Say that $H$ goes from $A$ to $B$ in $t$ steps within radius $r$ if the run of $H$ from
$A$ reaches $B$ after exactly $t$ steps, is in no exit state at any time before $t$, and at
every time up to $t$ keeps every work-tape head within distance $r$ of its position in $d$. Let
$i$ be a phase other than the stop phase with machine $M_i$, let $s,v$ be controls of phase
$i$, and let $q$ be a state of $H$ such that the redirect table sends $(i,v)$ to $q$ and
redirects no other control of phase $i$. Let $a,a',b,b'\in\mathbb{Z}$, $n,t,r\in\mathbb{N}$, a
tape $k_0$, $w\in\mathbb{Z}$ and $\ell,u\in\mathbb{Z}^{K+2}$. Suppose $M_i$ run from
$\omega(s,a,n,b)$ for $t$ steps yields $\omega(s,a,n,b)$ with state $v$ and the head of tape
$k_0$ at $w$, and that this configuration equals $\omega(v,a',n,b')$; that $M_i$ is not in
state $v$ before time $t$; that at every time up to $t$ the head of tape $k_0$ lies in
$[\ell_{k_0},u_{k_0}]$ and every other head is at its position in $\omega(s,a,n,b)$; that
$h_{k_0}-r\le\ell_{k_0}$ and $u_{k_0}\le h_{k_0}+r$; and that $\abs{a}\le r$ and $\abs{b}\le
r$. Then $H$ goes from $\omega((i,s),a,n,b)$ to $\omega(q,a',n,b')$ in $t+1$ steps within
radius $r$.
\end{lemma}

\begin{lemma}[Copying a field onto the clock]
\lean{Turing.cw_seg_copy}
\label{Turing.cw_seg_copy}
\leanok
\uses{Turing.Cfg, Turing.FinTM.bufferTape, Turing.MultiTapeTM, Turing.MultiTapeTM.runFrom, Turing.SeekPhase, Turing.Segment, Turing.SweepPhase, Turing.WindowExit, Turing.copyMapTM, Turing.cw_Ctrl, Turing.cw_Idx, Turing.cw_box, Turing.cw_box_of, Turing.cw_cfg, Turing.cw_cfg_clockTape, Turing.cw_cfg_mapState, Turing.cw_cfg_pos, Turing.cw_cfg_posL, Turing.cw_cfg_pos₀, Turing.cw_cfg_tape, Turing.cw_copy_constant, Turing.cw_guard, Turing.cw_hostTM, Turing.cw_join, Turing.cw_red, Turing.cw_seg, Turing.cw_zero_ne_last}
\difficulty{5}
Let $K,c,m\in\mathbb{N}$, let $e$ be a tape symbol, and let $R$ be a Turing machine with $K+2$
work tapes whose state set $P$ has decidable equality, together with an entry state
$p_{\mathrm{en}}\in P$ and a map $\mathrm{ex}$ from the three window exits (halted, left,
exhausted) to $P$. Let $H$ be the clock runner built from these data: the phase host on $K+2$
work tapes whose phases are the clock writer's phases for $c$ passes over $m$ record fields
(acting on the record tape $0$ and the clock tape $K+1$), the runner $R$, the clear of the
clock tape, and a stationary stop phase; the states of the stop phase are called exits. For a
configuration $d$ of $H$, a control $q$, integers $a,b$ and $n\in\mathbb{N}$, let
$\omega(q,a,n,b)$ be the configuration obtained from $d$ by setting the state to $q$, moving
the head of tape $0$ by $a$ and the head of tape $K+1$ by $b$, and writing $n$ copies of $e$ on
tape $K+1$ starting at its head position in $d$; let $h_j$ denote the head position of tape $j$
in $d$. Say that $H$ goes from $A$ to $B$ in $t$ steps within radius $r$ if the run of $H$ from
$A$ reaches $B$ after exactly $t$ steps, is in no exit state at any time before $t$, and at
every time up to $t$ keeps every work-tape head within distance $r$ of its position in $d$. Let
$p<c$, $j<m$, let $f$ be a word, $a\in\mathbb{Z}$ and $n,n',r\in\mathbb{N}$. Suppose that for
every integer $-1\le z\le|f|$ the cell $h_0+a+z$ of tape $0$ in $d$ holds the $z$-th letter of
$f$ (blank at $z=-1$ and $z=|f|$), that $n'=n+|f|$, $0\le a$, $a+|f|\le r$ and $n'+1\le r$.
Then $H$ goes from $\omega(q_1,a,n,n)$ to $\omega(q_2,a,n',n)$ in $2|f|+3$ steps within radius
$r$, where $q_1$ is the copy phase of field $j$ in pass $p$ with control sweep and $q_2$ is the
record skip of $(p,j)$ with control run.
\end{lemma}

\begin{lemma}[Skipping a field on the record tape]
\lean{Turing.cw_seg_skipS}
\label{Turing.cw_seg_skipS}
\leanok
\uses{Turing.Cfg, Turing.FinTM.bufferTape, Turing.MultiTapeTM, Turing.SeekPhase, Turing.Segment, Turing.WindowExit, Turing.cw_Ctrl, Turing.cw_Idx, Turing.cw_box, Turing.cw_cfg, Turing.cw_cfg_move₀, Turing.cw_cfg_pos₀, Turing.cw_cfg_tape, Turing.cw_hostTM, Turing.cw_seg_row, Turing.cw_zero_ne_last, Turing.skipTM_run_ofCfg}
\difficulty{3}
Let $K,c,m\in\mathbb{N}$, let $e$ be a tape symbol, and let $R$ be a Turing machine with $K+2$
work tapes whose state set $P$ has decidable equality, together with an entry state
$p_{\mathrm{en}}\in P$ and a map $\mathrm{ex}$ from the three window exits (halted, left,
exhausted) to $P$. Let $H$ be the clock runner built from these data: the phase host on $K+2$
work tapes whose phases are the clock writer's phases for $c$ passes over $m$ record fields
(acting on the record tape $0$ and the clock tape $K+1$), the runner $R$, the clear of the
clock tape, and a stationary stop phase; the states of the stop phase are called exits. For a
configuration $d$ of $H$, a control $q$, integers $a,b$ and $n\in\mathbb{N}$, let
$\omega(q,a,n,b)$ be the configuration obtained from $d$ by setting the state to $q$, moving
the head of tape $0$ by $a$ and the head of tape $K+1$ by $b$, and writing $n$ copies of $e$ on
tape $K+1$ starting at its head position in $d$; let $h_j$ denote the head position of tape $j$
in $d$. Say that $H$ goes from $A$ to $B$ in $t$ steps within radius $r$ if the run of $H$ from
$A$ reaches $B$ after exactly $t$ steps, is in no exit state at any time before $t$, and at
every time up to $t$ keeps every work-tape head within distance $r$ of its position in $d$. Let
$p<c$, $j<m$, let $f$ be a word, $a,a',b\in\mathbb{Z}$ and $n,r\in\mathbb{N}$. Suppose that for
every integer $0\le z\le|f|$ the cell $h_0+a+z$ of tape $0$ in $d$ holds the $z$-th letter of
$f$ (blank at $z=|f|$), that $a'=a+|f|+1$, $0\le a$, $a'\le r$ and $\abs{b}\le r$. Then $H$
goes from $\omega(q_1,a,n,b)$ to $\omega(q_2,a',n,b)$ in $|f|+2$ steps within radius $r$, where
$q_1$ is the record skip of $(p,j)$ with control run and $q_2$ is the clock skip indexed by
$(p,j)$ with control run.
\end{lemma}

\begin{lemma}[The word under the clock head]
\lean{Turing.cw_skip_word}
\label{Turing.cw_skip_word}
\leanok
\uses{Turing.Cfg, Turing.FinTM.bufferTape, Turing.cw_buffer_rep, Turing.cw_cfg, Turing.cw_cfg_clock, Turing.cw_cfg_posL}
\difficulty{4}
Let $K\in\mathbb{N}$, let $d$ be a configuration of a Turing machine with $K+2$ work tapes on
some input, and let $e$ be a tape symbol. For a state $q$, integers $a,b$ and $n\in\mathbb{N}$,
let $\omega(q,a,n,b)$ be the configuration obtained from $d$ by setting the state to $q$,
moving the head of tape $0$ by $a$ and the head of tape $K+1$ by $b$, and writing $n$ copies of
$e$ on tape $K+1$ starting at its head position in $d$; let $h_j$ denote the head position of
tape $j$ in $d$. Suppose tape $K+1$ of $d$ is entirely blank, and let $N,b,l\in\mathbb{N}$ with
$N=b+l$. Then for every integer $z\le l$ with either $z\ge 0$, or $b=0$ and $z\ge -1$, the cell
of tape $K+1$ in $\omega(q,a,N,b)$ at offset $z$ from its head holds $e$ if $0\le z<l$ and is
blank otherwise.
\end{lemma}

\begin{lemma}[Skipping the written clock cells]
\lean{Turing.cw_seg_skipC}
\label{Turing.cw_seg_skipC}
\leanok
\uses{Turing.Cfg, Turing.MultiTapeTM, Turing.SeekPhase, Turing.Segment, Turing.WindowExit, Turing.cw_Ctrl, Turing.cw_Idx, Turing.cw_afterSkip, Turing.cw_box, Turing.cw_cfg, Turing.cw_cfg_moveL, Turing.cw_cfg_posL, Turing.cw_hostTM, Turing.cw_seg_row, Turing.cw_skip_word, Turing.skipTM_run_ofCfg}
\difficulty{3}
Let $K,c,m\in\mathbb{N}$, let $e$ be a tape symbol, and let $R$ be a Turing machine with $K+2$
work tapes whose state set $P$ has decidable equality, together with an entry state
$p_{\mathrm{en}}\in P$ and a map $\mathrm{ex}$ from the three window exits (halted, left,
exhausted) to $P$. Let $H$ be the clock runner built from these data: the phase host on $K+2$
work tapes whose phases are the clock writer's phases for $c$ passes over $m$ record fields
(acting on the record tape $0$ and the clock tape $K+1$), the runner $R$, the clear of the
clock tape, and a stationary stop phase; the states of the stop phase are called exits. For a
configuration $d$ of $H$, a control $q$, integers $a,b$ and $n\in\mathbb{N}$, let
$\omega(q,a,n,b)$ be the configuration obtained from $d$ by setting the state to $q$, moving
the head of tape $0$ by $a$ and the head of tape $K+1$ by $b$, and writing $n$ copies of $e$ on
tape $K+1$ starting at its head position in $d$; let $h_j$ denote the head position of tape $j$
in $d$. Say that $H$ goes from $A$ to $B$ in $t$ steps within radius $r$ if the run of $H$ from
$A$ reaches $B$ after exactly $t$ steps, is in no exit state at any time before $t$, and at
every time up to $t$ keeps every work-tape head within distance $r$ of its position in $d$. Let
$o$ be either a pair $(p,j)$ with $p<c$, $j<m$, or the marker for the end of the last pass.
Suppose tape $K+1$ of $d$ is entirely blank, and let $a\in\mathbb{Z}$, $N,b,l,r\in\mathbb{N}$
and $b'\in\mathbb{Z}$ with $N=b+l$, $b'=b+l+1$, $\abs{a}\le r$ and $b'\le r$. Then $H$ goes
from $\omega(q_o,a,N,b)$ to $\omega(\sigma(o),a,N,b')$ in $l+2$ steps within radius $r$, where
$q_o$ is the clock skip indexed by $o$ with control run, $\sigma(p,j)$ is the left move of
$(p,j)$ with control false, and $\sigma$ of the end marker is the clock skip-back with control
run.
\end{lemma}

\begin{lemma}[Moving back onto the clock delimiter]
\lean{Turing.cw_seg_left}
\label{Turing.cw_seg_left}
\leanok
\uses{Turing.Cfg, Turing.MultiTapeTM, Turing.MultiTapeTM.runFrom, Turing.MultiTapeTM.runFrom_zero, Turing.Segment, Turing.WindowExit, Turing.cw_Ctrl, Turing.cw_Idx, Turing.cw_box, Turing.cw_cfg, Turing.cw_cfg_box, Turing.cw_cfg_mapState, Turing.cw_cfg_moveL, Turing.cw_cfg_posL, Turing.cw_guard, Turing.cw_hostTM, Turing.cw_leftTM, Turing.cw_left_run, Turing.cw_next, Turing.cw_red, Turing.cw_seg}
\difficulty{4}
Let $K,c,m\in\mathbb{N}$, let $e$ be a tape symbol, and let $R$ be a Turing machine with $K+2$
work tapes whose state set $P$ has decidable equality, together with an entry state
$p_{\mathrm{en}}\in P$ and a map $\mathrm{ex}$ from the three window exits (halted, left,
exhausted) to $P$. Let $H$ be the clock runner built from these data: the phase host on $K+2$
work tapes whose phases are the clock writer's phases for $c$ passes over $m$ record fields
(acting on the record tape $0$ and the clock tape $K+1$), the runner $R$, the clear of the
clock tape, and a stationary stop phase; the states of the stop phase are called exits. For a
configuration $d$ of $H$, a control $q$, integers $a,b$ and $n\in\mathbb{N}$, let
$\omega(q,a,n,b)$ be the configuration obtained from $d$ by setting the state to $q$, moving
the head of tape $0$ by $a$ and the head of tape $K+1$ by $b$, and writing $n$ copies of $e$ on
tape $K+1$ starting at its head position in $d$; let $h_j$ denote the head position of tape $j$
in $d$. Say that $H$ goes from $A$ to $B$ in $t$ steps within radius $r$ if the run of $H$ from
$A$ reaches $B$ after exactly $t$ steps, is in no exit state at any time before $t$, and at
every time up to $t$ keeps every work-tape head within distance $r$ of its position in $d$. Let
$p<c$, $j<m$, $a,b,b'\in\mathbb{Z}$ and $n,r\in\mathbb{N}$ with $b'=b-1$, $\abs{a}\le r$ and
$1\le b\le r$. Then $H$ goes from $\omega(q,a,n,b)$ to $\omega(\nu(p,j),a,n,b')$ in $2$ steps
within radius $r$, where $q$ is the left move of $(p,j)$ with control false and $\nu(p,j)$ is
the copy phase of field $j+1$ in pass $p$ with control sweep if $j+1<m$, and otherwise the skip
back over field $j$ in pass $p$ with control run.
\end{lemma}

\begin{lemma}[Restating a segment's length]
\lean{Turing.cw_seg_time}
\label{Turing.cw_seg_time}
\leanok
\uses{Turing.Cfg, Turing.MultiTapeTM, Turing.Segment}
\difficulty{1}
Let $M$ be a multi-tape Turing machine, $F$ a set of states, $B$ a property of head vectors,
$A,A'$ configurations and $t,t'\in\mathbb{N}$ with $t=t'$. If the run of $M$ from $A$ reaches
$A'$ after exactly $t$ steps, avoids $F$ before time $t$ and satisfies $B$ at every time up to
$t$, then the same holds with $t'$ in place of $t$.
\end{lemma}

\begin{lemma}[A record field as a delimited word]
\lean{Turing.cw_field}
\label{Turing.cw_field}
\leanok
\uses{Turing.Cfg, Turing.FinTM.bufferTape, Turing.RecordAt, Turing.fieldOffset, Turing.fieldOffset_add_length_le, Turing.fieldsTape_field, Turing.recordAt_iff}
\difficulty{3}
Let $K\in\mathbb{N}$, let $d$ be a configuration with $K$ work tapes, let $r$ be a work tape,
and let $f_0,\dots,f_{\ell-1}$ be words such that tape $r$ holds, from its head position $h$,
the fields $f_0,\dots,f_{\ell-1}$ separated by single blanks and delimited by blanks at offsets
$-1$ and just after the layout. Let $i<\ell$ and $o_i=\sum_{j<i}(|f_j|+1)$. Then for every
integer $-1\le z\le|f_i|$, the cell $h+o_i+z$ of tape $r$ holds the $z$-th letter of $f_i$, and
is blank at $z=-1$ and at $z=|f_i|$.
\end{lemma}

\begin{lemma}[One field's round of the clock writer]
\lean{Turing.cw_round}
\label{Turing.cw_round}
\leanok
\uses{Turing.Cfg, Turing.MultiTapeTM, Turing.RecordAt, Turing.Segment, Turing.Segment.trans, Turing.SweepPhase, Turing.WindowExit, Turing.cw_Ctrl, Turing.cw_Idx, Turing.cw_box, Turing.cw_cfg, Turing.cw_field, Turing.cw_hostTM, Turing.cw_next, Turing.cw_seg_copy, Turing.cw_seg_left, Turing.cw_seg_skipC, Turing.cw_seg_skipS, Turing.cw_seg_time, Turing.fieldOffset, Turing.fieldOffset_succ}
\difficulty{5}
Let $K,c,m\in\mathbb{N}$, let $e$ be a tape symbol, and let $R$ be a Turing machine with $K+2$
work tapes whose state set $P$ has decidable equality, together with an entry state
$p_{\mathrm{en}}\in P$ and a map $\mathrm{ex}$ from the three window exits (halted, left,
exhausted) to $P$. Let $H$ be the clock runner built from these data: the phase host on $K+2$
work tapes whose phases are the clock writer's phases for $c$ passes over $m$ record fields
(acting on the record tape $0$ and the clock tape $K+1$), the runner $R$, the clear of the
clock tape, and a stationary stop phase; the states of the stop phase are called exits. For a
configuration $d$ of $H$, a control $q$, integers $a,b$ and $n\in\mathbb{N}$, let
$\omega(q,a,n,b)$ be the configuration obtained from $d$ by setting the state to $q$, moving
the head of tape $0$ by $a$ and the head of tape $K+1$ by $b$, and writing $n$ copies of $e$ on
tape $K+1$ starting at its head position in $d$; let $h_j$ denote the head position of tape $j$
in $d$. Say that $H$ goes from $A$ to $B$ in $t$ steps within radius $r$ if the run of $H$ from
$A$ reaches $B$ after exactly $t$ steps, is in no exit state at any time before $t$, and at
every time up to $t$ keeps every work-tape head within distance $r$ of its position in $d$. Let
$p<c$, let $f_0,\dots,f_{m-1}$ be words, let $d$ be a configuration of $H$ whose tape $0$
holds, from its head, the fields $f_0,\dots,f_{m-1}$ separated by single blanks and delimited
by blanks on both sides, and whose tape $K+1$ is entirely blank, and write
$o_j=\sum_{i<j}(|f_i|+1)$ for the offset of field $j$. Let $j<m$ and $n,r\in\mathbb{N}$ with
$n+|f_j|+1\le r$ and $o_{j+1}\le r$. Then $H$ goes from $\omega(q,o_j,n,n)$ to
$\omega(\nu(p,j),o_{j+1},n+|f_j|,n+|f_j|)$ in $4|f_j|+9$ steps within radius $r$, where $q$ is
the copy phase of field $j$ in pass $p$ with control sweep and $\nu(p,j)$ is the copy phase of
field $j+1$ in pass $p$ with control sweep if $j+1<m$, and otherwise the skip back over field
$j$ in pass $p$ with control run.
\end{lemma}

\begin{lemma}[One pass of the clock writer]
\lean{Turing.cw_pass}
\label{Turing.cw_pass}
\leanok
\uses{Turing.Cfg, Turing.MultiTapeTM, Turing.RecordAt, Turing.SeekPhase, Turing.Segment, Turing.Segment.trans, Turing.SweepPhase, Turing.WindowExit, Turing.cw_Ctrl, Turing.cw_Idx, Turing.cw_box, Turing.cw_cfg, Turing.cw_hostTM, Turing.cw_next, Turing.cw_round, Turing.fieldOffset, Turing.fieldOffset_add_length_le, Turing.fieldOffset_succ, Turing.length_joinFields_add_one}
\difficulty{6}
Let $K,c,m\in\mathbb{N}$, let $e$ be a tape symbol, and let $R$ be a Turing machine with $K+2$
work tapes whose state set $P$ has decidable equality, together with an entry state
$p_{\mathrm{en}}\in P$ and a map $\mathrm{ex}$ from the three window exits (halted, left,
exhausted) to $P$. Let $H$ be the clock runner built from these data: the phase host on $K+2$
work tapes whose phases are the clock writer's phases for $c$ passes over $m$ record fields
(acting on the record tape $0$ and the clock tape $K+1$), the runner $R$, the clear of the
clock tape, and a stationary stop phase; the states of the stop phase are called exits. For a
configuration $d$ of $H$, a control $q$, integers $a,b$ and $n\in\mathbb{N}$, let
$\omega(q,a,n,b)$ be the configuration obtained from $d$ by setting the state to $q$, moving
the head of tape $0$ by $a$ and the head of tape $K+1$ by $b$, and writing $n$ copies of $e$ on
tape $K+1$ starting at its head position in $d$; let $h_j$ denote the head position of tape $j$
in $d$. Say that $H$ goes from $A$ to $B$ in $t$ steps within radius $r$ if the run of $H$ from
$A$ reaches $B$ after exactly $t$ steps, is in no exit state at any time before $t$, and at
every time up to $t$ keeps every work-tape head within distance $r$ of its position in $d$. Let
$p<c$, let $f_0,\dots,f_{m-1}$ be words with $m\ge 1$, let $d$ be a configuration of $H$ whose
tape $0$ holds, from its head, the fields $f_0,\dots,f_{m-1}$ separated by single blanks and
delimited by blanks on both sides, and whose tape $K+1$ is entirely blank, and write
$o_j=\sum_{i<j}(|f_i|+1)$ for the offset of field $j$. Let $N_0,r\in\mathbb{N}$ with
$N_0+o_m+1\le r$. Then for every $j<m$ and all $n,n'\in\mathbb{N}$ with $n+j=N_0+o_j$ and
$n'+m=N_0+o_m$, there is $t$ with $t+9o_j\le 9o_m$ such that $H$ goes from
$\omega(q_j,o_j,n,n)$ to $\omega(q',o_m,n',n')$ in $t$ steps within radius $r$, where $q_j$ is
the copy phase of field $j$ in pass $p$ with control sweep and $q'$ is the skip back over field
$m-1$ in pass $p$ with control run.
\end{lemma}

\begin{lemma}[Skipping back over a record field]
\lean{Turing.cw_seg_back}
\label{Turing.cw_seg_back}
\leanok
\uses{Turing.Cfg, Turing.FinTM.bufferTape, Turing.MultiTapeTM, Turing.SeekPhase, Turing.Segment, Turing.WindowExit, Turing.cw_Ctrl, Turing.cw_Idx, Turing.cw_box, Turing.cw_cfg, Turing.cw_cfg_move₀, Turing.cw_cfg_pos₀, Turing.cw_cfg_tape, Turing.cw_hostTM, Turing.cw_prev, Turing.cw_seg_row, Turing.cw_zero_ne_last, Turing.skipBackTM_run_ofCfg}
\difficulty{3}
Let $K,c,m\in\mathbb{N}$, let $e$ be a tape symbol, and let $R$ be a Turing machine with $K+2$
work tapes whose state set $P$ has decidable equality, together with an entry state
$p_{\mathrm{en}}\in P$ and a map $\mathrm{ex}$ from the three window exits (halted, left,
exhausted) to $P$. Let $H$ be the clock runner built from these data: the phase host on $K+2$
work tapes whose phases are the clock writer's phases for $c$ passes over $m$ record fields
(acting on the record tape $0$ and the clock tape $K+1$), the runner $R$, the clear of the
clock tape, and a stationary stop phase; the states of the stop phase are called exits. For a
configuration $d$ of $H$, a control $q$, integers $a,b$ and $n\in\mathbb{N}$, let
$\omega(q,a,n,b)$ be the configuration obtained from $d$ by setting the state to $q$, moving
the head of tape $0$ by $a$ and the head of tape $K+1$ by $b$, and writing $n$ copies of $e$ on
tape $K+1$ starting at its head position in $d$; let $h_j$ denote the head position of tape $j$
in $d$. Say that $H$ goes from $A$ to $B$ in $t$ steps within radius $r$ if the run of $H$ from
$A$ reaches $B$ after exactly $t$ steps, is in no exit state at any time before $t$, and at
every time up to $t$ keeps every work-tape head within distance $r$ of its position in $d$. Let
$p<c$, $j<m$, let $f$ be a word, $a,a',b\in\mathbb{Z}$ and $n,r\in\mathbb{N}$. Suppose that for
every integer $-1\le z\le|f|$ the cell $h_0+a'+z$ of tape $0$ in $d$ holds the $z$-th letter of
$f$ (blank at $z=-1$ and $z=|f|$), that $a'=a-|f|-1$, $0\le a'$, $a\le r$, $\abs{b}\le r$ and
$1\le r$. Then $H$ goes from $\omega(q,a,n,b)$ to $\omega(\pi(p,j),a',n,b)$ in $|f|+4$ steps
within radius $r$, where $q$ is the skip back over field $j$ in pass $p$ with control run and
$\pi(p,j)$ is the skip back over field $j-1$ in pass $p$ with control run if $j>0$; otherwise
the copy phase of field $0$ in pass $p+1$ with control sweep if $p+1<c$; and otherwise the
final clock skip (after the last pass) with control run.
\end{lemma}

\begin{lemma}[Rewinding the record tape]
\lean{Turing.cw_rewind}
\label{Turing.cw_rewind}
\leanok
\uses{Turing.Cfg, Turing.MultiTapeTM, Turing.RecordAt, Turing.SeekPhase, Turing.Segment, Turing.Segment.trans, Turing.WindowExit, Turing.cw_Ctrl, Turing.cw_Idx, Turing.cw_box, Turing.cw_cfg, Turing.cw_field, Turing.cw_hostTM, Turing.cw_prev, Turing.cw_seg_back, Turing.fieldOffset, Turing.fieldOffset_add_length_le, Turing.fieldOffset_succ, Turing.fieldOffset_zero, Turing.length_joinFields_add_one}
\difficulty{5}
Let $K,c,m\in\mathbb{N}$, let $e$ be a tape symbol, and let $R$ be a Turing machine with $K+2$
work tapes whose state set $P$ has decidable equality, together with an entry state
$p_{\mathrm{en}}\in P$ and a map $\mathrm{ex}$ from the three window exits (halted, left,
exhausted) to $P$. Let $H$ be the clock runner built from these data: the phase host on $K+2$
work tapes whose phases are the clock writer's phases for $c$ passes over $m$ record fields
(acting on the record tape $0$ and the clock tape $K+1$), the runner $R$, the clear of the
clock tape, and a stationary stop phase; the states of the stop phase are called exits. For a
configuration $d$ of $H$, a control $q$, integers $a,b$ and $n\in\mathbb{N}$, let
$\omega(q,a,n,b)$ be the configuration obtained from $d$ by setting the state to $q$, moving
the head of tape $0$ by $a$ and the head of tape $K+1$ by $b$, and writing $n$ copies of $e$ on
tape $K+1$ starting at its head position in $d$; let $h_j$ denote the head position of tape $j$
in $d$. Say that $H$ goes from $A$ to $B$ in $t$ steps within radius $r$ if the run of $H$ from
$A$ reaches $B$ after exactly $t$ steps, is in no exit state at any time before $t$, and at
every time up to $t$ keeps every work-tape head within distance $r$ of its position in $d$. Let
$p<c$, let $f_0,\dots,f_{m-1}$ be words with $m>0$, let $d$ be a configuration of $H$ whose
tape $0$ holds, from its head, the fields $f_0,\dots,f_{m-1}$ separated by single blanks and
delimited by blanks on both sides, and write $o_j=\sum_{i<j}(|f_i|+1)$ for the offset of field
$j$. Let $n,r\in\mathbb{N}$ and $b\in\mathbb{Z}$ with $o_m\le r$ and $\abs{b}\le r$. Then for
every $j<m$ there is $t\le 4o_{j+1}$ such that $H$ goes from $\omega(q_j,o_{j+1},n,b)$ to
$\omega(\pi(p,0),0,n,b)$ in $t$ steps within radius $r$, where $q_j$ is the skip back over
field $j$ in pass $p$ with control run and $\pi(p,j)$ is the skip back over field $j-1$ in pass
$p$ with control run if $j>0$; otherwise the copy phase of field $0$ in pass $p+1$ with control
sweep if $p+1<c$; and otherwise the final clock skip (after the last pass) with control run.
\end{lemma}

\begin{lemma}[A pass followed by the record rewind]
\lean{Turing.cw_pass_rewind}
\label{Turing.cw_pass_rewind}
\leanok
\uses{Turing.Cfg, Turing.MultiTapeTM, Turing.RecordAt, Turing.Segment, Turing.Segment.trans, Turing.SweepPhase, Turing.WindowExit, Turing.cw_Ctrl, Turing.cw_Idx, Turing.cw_box, Turing.cw_cfg, Turing.cw_hostTM, Turing.cw_pass, Turing.cw_prev, Turing.cw_rewind, Turing.fieldOffset, Turing.fieldOffset_zero}
\difficulty{4}
Let $K,c,m\in\mathbb{N}$, let $e$ be a tape symbol, and let $R$ be a Turing machine with $K+2$
work tapes whose state set $P$ has decidable equality, together with an entry state
$p_{\mathrm{en}}\in P$ and a map $\mathrm{ex}$ from the three window exits (halted, left,
exhausted) to $P$. Let $H$ be the clock runner built from these data: the phase host on $K+2$
work tapes whose phases are the clock writer's phases for $c$ passes over $m$ record fields
(acting on the record tape $0$ and the clock tape $K+1$), the runner $R$, the clear of the
clock tape, and a stationary stop phase; the states of the stop phase are called exits. For a
configuration $d$ of $H$, a control $q$, integers $a,b$ and $n\in\mathbb{N}$, let
$\omega(q,a,n,b)$ be the configuration obtained from $d$ by setting the state to $q$, moving
the head of tape $0$ by $a$ and the head of tape $K+1$ by $b$, and writing $n$ copies of $e$ on
tape $K+1$ starting at its head position in $d$; let $h_j$ denote the head position of tape $j$
in $d$. Say that $H$ goes from $A$ to $B$ in $t$ steps within radius $r$ if the run of $H$ from
$A$ reaches $B$ after exactly $t$ steps, is in no exit state at any time before $t$, and at
every time up to $t$ keeps every work-tape head within distance $r$ of its position in $d$. Let
$p<c$, let $f_0,\dots,f_{m-1}$ be words with $m>0$, let $d$ be a configuration of $H$ whose
tape $0$ holds, from its head, the fields $f_0,\dots,f_{m-1}$ separated by single blanks and
delimited by blanks on both sides, and whose tape $K+1$ is entirely blank, and write
$o_j=\sum_{i<j}(|f_i|+1)$ for the offset of field $j$. Let $F,N_0,N_1,r\in\mathbb{N}$ with
$o_m=F+m$, $N_1=N_0+F$ and $N_0+o_m+1\le r$. Then there is $t\le 13o_m$ such that $H$ goes from
$\omega(q,0,N_0,N_0)$ to $\omega(\pi(p,0),0,N_1,N_1)$ in $t$ steps within radius $r$, where $q$
is the copy phase of field $0$ in pass $p$ with control sweep and $\pi(p,j)$ is the skip back
over field $j-1$ in pass $p$ with control run if $j>0$; otherwise the copy phase of field $0$
in pass $p+1$ with control sweep if $p+1<c$; and otherwise the final clock skip (after the last
pass) with control run.
\end{lemma}

\begin{lemma}[All passes of the clock writer]
\lean{Turing.cw_passes}
\label{Turing.cw_passes}
\leanok
\uses{Turing.Cfg, Turing.MultiTapeTM, Turing.RecordAt, Turing.SeekPhase, Turing.Segment, Turing.Segment.trans, Turing.SweepPhase, Turing.WindowExit, Turing.cw_Ctrl, Turing.cw_Idx, Turing.cw_box, Turing.cw_cfg, Turing.cw_hostTM, Turing.cw_pass_rewind, Turing.cw_prev, Turing.fieldOffset}
\difficulty{6}
Let $K,c,m\in\mathbb{N}$, let $e$ be a tape symbol, and let $R$ be a Turing machine with $K+2$
work tapes whose state set $P$ has decidable equality, together with an entry state
$p_{\mathrm{en}}\in P$ and a map $\mathrm{ex}$ from the three window exits (halted, left,
exhausted) to $P$. Let $H$ be the clock runner built from these data: the phase host on $K+2$
work tapes whose phases are the clock writer's phases for $c$ passes over $m$ record fields
(acting on the record tape $0$ and the clock tape $K+1$), the runner $R$, the clear of the
clock tape, and a stationary stop phase; the states of the stop phase are called exits. For a
configuration $d$ of $H$, a control $q$, integers $a,b$ and $n\in\mathbb{N}$, let
$\omega(q,a,n,b)$ be the configuration obtained from $d$ by setting the state to $q$, moving
the head of tape $0$ by $a$ and the head of tape $K+1$ by $b$, and writing $n$ copies of $e$ on
tape $K+1$ starting at its head position in $d$; let $h_j$ denote the head position of tape $j$
in $d$. Say that $H$ goes from $A$ to $B$ in $t$ steps within radius $r$ if the run of $H$ from
$A$ reaches $B$ after exactly $t$ steps, is in no exit state at any time before $t$, and at
every time up to $t$ keeps every work-tape head within distance $r$ of its position in $d$. Let
$f_0,\dots,f_{m-1}$ be words with $m>0$, let $d$ be a configuration of $H$ whose tape $0$
holds, from its head, the fields $f_0,\dots,f_{m-1}$ separated by single blanks and delimited
by blanks on both sides, and whose tape $K+1$ is entirely blank, and write
$o_j=\sum_{i<j}(|f_i|+1)$ for the offset of field $j$. Let $F,r\in\mathbb{N}$ with $o_m=F+m$
and $c\,o_m+1\le r$. Then for every pass $p<c$ there is $t$ with $t+13p\,o_m\le 13c\,o_m$ such
that $H$ goes from $\omega(q_p,0,pF,pF)$ to $\omega(q',0,cF,cF)$ in $t$ steps within radius
$r$, where $q_p$ is the copy phase of field $0$ in pass $p$ with control sweep and $q'$ is the
final clock skip (after the last pass) with control run.
\end{lemma}

\begin{lemma}[Skipping back over the clock word]
\lean{Turing.cw_seg_clockBack}
\label{Turing.cw_seg_clockBack}
\leanok
\uses{Turing.Cfg, Turing.FinTM.bufferTape, Turing.MultiTapeTM, Turing.SeekPhase, Turing.Segment, Turing.WindowExit, Turing.cw_Ctrl, Turing.cw_Idx, Turing.cw_box, Turing.cw_cfg, Turing.cw_cfg_clockTape, Turing.cw_cfg_moveL, Turing.cw_cfg_posL, Turing.cw_hostTM, Turing.cw_seg_row, Turing.cw_skip_word, Turing.skipBackTM_run_ofCfg}
\difficulty{4}
Let $K,c,m\in\mathbb{N}$, let $e$ be a tape symbol, and let $R$ be a Turing machine with $K+2$
work tapes whose state set $P$ has decidable equality, together with an entry state
$p_{\mathrm{en}}\in P$ and a map $\mathrm{ex}$ from the three window exits (halted, left,
exhausted) to $P$. Let $H$ be the clock runner built from these data: the phase host on $K+2$
work tapes whose phases are the clock writer's phases for $c$ passes over $m$ record fields
(acting on the record tape $0$ and the clock tape $K+1$), the runner $R$, the clear of the
clock tape, and a stationary stop phase; the states of the stop phase are called exits. For a
configuration $d$ of $H$, a control $q$, integers $a,b$ and $n\in\mathbb{N}$, let
$\omega(q,a,n,b)$ be the configuration obtained from $d$ by setting the state to $q$, moving
the head of tape $0$ by $a$ and the head of tape $K+1$ by $b$, and writing $n$ copies of $e$ on
tape $K+1$ starting at its head position in $d$; let $h_j$ denote the head position of tape $j$
in $d$. Say that $H$ goes from $A$ to $B$ in $t$ steps within radius $r$ if the run of $H$ from
$A$ reaches $B$ after exactly $t$ steps, is in no exit state at any time before $t$, and at
every time up to $t$ keeps every work-tape head within distance $r$ of its position in $d$.
Suppose tape $K+1$ of $d$ is entirely blank, and let $W,r\in\mathbb{N}$ with $W+1\le r$. Then
$H$ goes from $\omega(q,0,W,W+1)$ to $\omega(q',0,W,0)$ in $W+4$ steps within radius $r$, where
$q$ is the clock skip-back with control run and $q'$ is the runner phase at $p_{\mathrm{en}}$.
\end{lemma}

\begin{lemma}[Unchanged final output means silent run]
\lean{Turing.cw_silence}
\label{Turing.cw_silence}
\leanok
\uses{Turing.Cfg, Turing.MultiTapeTM, Turing.MultiTapeTM.output_prefix, Turing.MultiTapeTM.runFrom, Turing.MultiTapeTM.runFrom_zero}
\difficulty{4}
Let $M$ be a multi-tape Turing machine, $d$ a configuration and $t\in\mathbb{N}$. If the output
of the run of $M$ from $d$ after $t$ steps equals the output of $d$, then at every time $u\le
t$ the output of the run equals the output of $d$.
\end{lemma}

\begin{lemma}[The clock writer's contract]
\lean{Turing.cw_writer}
\label{Turing.cw_writer}
\leanok
\uses{Turing.Cfg, Turing.MultiTapeTM, Turing.MultiTapeTM.runFrom, Turing.RecordAt, Turing.Segment, Turing.Segment.trans, Turing.SweepPhase, Turing.WindowExit, Turing.cw_Ctrl, Turing.cw_Idx, Turing.cw_box, Turing.cw_cfg, Turing.cw_cfg_zero, Turing.cw_hostTM, Turing.cw_passes, Turing.cw_seg_clockBack, Turing.cw_seg_skipC, Turing.cw_silence, Turing.fieldOffset, Turing.joinFields_length, Turing.length_joinFields_add_one, Turing.segment_iff}
\difficulty{5}
Let $K,c,m\in\mathbb{N}$, let $e$ be a tape symbol, and let $R$ be a Turing machine with $K+2$
work tapes whose state set $P$ has decidable equality, together with an entry state
$p_{\mathrm{en}}\in P$ and a map $\mathrm{ex}$ from the three window exits (halted, left,
exhausted) to $P$. Let $H$ be the clock runner built from these data: the phase host on $K+2$
work tapes whose phases are the clock writer's phases for $c$ passes over $m$ record fields
(acting on the record tape $0$ and the clock tape $K+1$), the runner $R$, the clear of the
clock tape, and a stationary stop phase; the states of the stop phase are called exits. For a
configuration $d$ of $H$, a control $q$, integers $a,b$ and $n\in\mathbb{N}$, let
$\omega(q,a,n,b)$ be the configuration obtained from $d$ by setting the state to $q$, moving
the head of tape $0$ by $a$ and the head of tape $K+1$ by $b$, and writing $n$ copies of $e$ on
tape $K+1$ starting at its head position in $d$; let $h_j$ denote the head position of tape $j$
in $d$. Say that $H$ goes from $A$ to $B$ in $t$ steps within radius $r$ if the run of $H$ from
$A$ reaches $B$ after exactly $t$ steps, is in no exit state at any time before $t$, and at
every time up to $t$ keeps every work-tape head within distance $r$ of its position in $d$. Let
$c>0$ and $m>0$, let $f_0,\dots,f_{m-1}$ be words, write $o_j=\sum_{i<j}(|f_i|+1)$ for the
offset of field $j$, and let $F=\sum_{j<m}|f_j|$. Let $d$ be a configuration of $H$ in the copy
phase of field $0$ in pass $0$ with control sweep, whose tape $0$ holds, from its head, the
fields $f_0,\dots,f_{m-1}$ separated by single blanks and delimited by blanks on both sides,
and whose tape $K+1$ is entirely blank, and let $r\in\mathbb{N}$ with $c\,o_m+2\le r$. Then
there is $t\le 13c\,o_m+cF+6$ such that $H$ goes from $d$ to $\omega(q',0,cF,0)$ in $t$ steps
within radius $r$, where $q'$ is the runner phase at $p_{\mathrm{en}}$, and the output of $H$
at every time up to $t$ equals that of $d$.
\end{lemma}

\begin{lemma}[The windowed runner as a host segment]
\lean{Turing.cw_seg_run}
\label{Turing.cw_seg_run}
\leanok
\uses{Turing.Cfg, Turing.Cfg.mapState, Turing.MultiTapeTM, Turing.MultiTapeTM.runFrom, Turing.Segment, Turing.WindowExit, Turing.cw_Ctrl, Turing.cw_Idx, Turing.cw_hostTM, Turing.cw_red, Turing.cw_seg}
\difficulty{3}
Let $K,c,m\in\mathbb{N}$, let $e$ be a tape symbol, and let $R$ be a Turing machine with $K+2$
work tapes whose state set $P$ has decidable equality, together with an entry state
$p_{\mathrm{en}}\in P$ and a map $\mathrm{ex}$ from the three window exits (halted, left,
exhausted) to $P$. Let $H$ be the clock runner built from these data: the phase host on $K+2$
work tapes whose phases are the clock writer's phases for $c$ passes over $m$ record fields
(acting on the record tape $0$ and the clock tape $K+1$), the runner $R$, the clear of the
clock tape, and a stationary stop phase; the states of the stop phase are called exits. Let $a$
be a configuration of $R$, $t\in\mathbb{N}$, $y$ a window exit and $q$ a state of $H$. Suppose
the run of $R$ from $a$ is in state $\mathrm{ex}(y)$ at time $t$ and in no state
$\mathrm{ex}(y')$ before time $t$, that the redirect table sends the runner phase at
$\mathrm{ex}(y)$ to $q$, and that $B$ is a property of head vectors holding at every time up to
$t$ of the run of $R$. Then the run of $H$ from $a$ with its state tagged by the runner phase
reaches the configuration of $R$ at time $t$ with its state replaced by $q$ after exactly $t+1$
steps, is in no exit state before time $t+1$, and its head vector satisfies $B$ at every time
up to $t+1$.
\end{lemma}

\begin{lemma}[Clearing the clock after the runner]
\lean{Turing.cw_seg_clear}
\label{Turing.cw_seg_clear}
\leanok
\uses{Turing.Cfg, Turing.Cfg.mapState, Turing.MultiTapeTM, Turing.Segment, Turing.SweepPhase, Turing.WindowExit, Turing.cw_Ctrl, Turing.cw_Idx, Turing.cw_box, Turing.cw_box_row, Turing.cw_cfg, Turing.cw_cfg_clock, Turing.cw_cfg_posL, Turing.cw_cfg_pos_zero, Turing.cw_cfg_pos₀, Turing.cw_cfg_tape, Turing.cw_clear_patch, Turing.cw_guard, Turing.cw_hostTM, Turing.cw_mapState_io, Turing.cw_mapState_pos, Turing.cw_mapState_tapes, Turing.cw_red, Turing.cw_seg, Turing.cw_zero_ne_last, Turing.wordPatch}
\difficulty{6}
Let $K,c,m\in\mathbb{N}$, let $e$ be a tape symbol, and let $R$ be a Turing machine with $K+2$
work tapes whose state set $P$ has decidable equality, together with an entry state
$p_{\mathrm{en}}\in P$ and a map $\mathrm{ex}$ from the three window exits (halted, left,
exhausted) to $P$. Let $H$ be the clock runner built from these data: the phase host on $K+2$
work tapes whose phases are the clock writer's phases for $c$ passes over $m$ record fields
(acting on the record tape $0$ and the clock tape $K+1$), the runner $R$, the clear of the
clock tape, and a stationary stop phase; the states of the stop phase are called exits. For a
configuration $d$ of $H$, a control $q$, integers $a,b$ and $n\in\mathbb{N}$, let
$\omega(q,a,n,b)$ be the configuration obtained from $d$ by setting the state to $q$, moving
the head of tape $0$ by $a$ and the head of tape $K+1$ by $b$, and writing $n$ copies of $e$ on
tape $K+1$ starting at its head position in $d$; let $h_j$ denote the head position of tape $j$
in $d$. Say that $H$ goes from $A$ to $B$ in $t$ steps within radius $r$ if the run of $H$ from
$A$ reaches $B$ after exactly $t$ steps, is in no exit state at any time before $t$, and at
every time up to $t$ keeps every work-tape head within distance $r$ of its position in $d$. Let
$b$ be a Boolean, let $d$ be a configuration of $H$ whose tape $K+1$ is entirely blank, let
$W,r\in\mathbb{N}$ with $W+1\le r$, let $w$ be a word with $|w|=W$, let $L\in\mathbb{Z}$, let
$g$ assign a cell content to each integer, and let $y\in P$. Let $A$ be
$\omega(p_{\mathrm{en}},0,W,0)$ with state $y$, with each cell $z$ of tape $0$ satisfying
$h_0\le z<h_0+L$ replaced by $g(z-h_0)$, and with $w$ written on tape $K+1$ from $h_{K+1}$;
then retag its state as the clear phase tagged $b$ with control sweep. Let $B$ be $d$ with the
same replacement on tape $0$ and with state the stop phase at the exit halted if $b$ is true
and left otherwise. Then $H$ goes from $A$ to $B$ in $2W+3$ steps within radius $r$.
\end{lemma}

\begin{lemma}[Run-time clauses from a segment]
\lean{Turing.cw_finish}
\label{Turing.cw_finish}
\leanok
\uses{Turing.Cfg, Turing.MultiTapeTM, Turing.MultiTapeTM.runFrom, Turing.MultiTapeTM.spaceUsed, Turing.MultiTapeTM.spaceUsed_le_of_abs_le, Turing.Segment, Turing.cw_box, Turing.cw_silence, Turing.segment_iff}
\difficulty{2}
Let $H$ be a Turing machine with $k$ work tapes, $F$ a set of states, $d,b$ configurations and
$t,r\in\mathbb{N}$. Suppose the run of $H$ from $d$ reaches $b$ after exactly $t$ steps, is in
no state of $F$ before time $t$, keeps every work-tape head within distance $r$ of its position
in $d$ at every time up to $t$, and that $b$ has the same output as $d$. Then the run is in no
state of $F$ before time $t$, every head stays within distance $r$ of its initial position up
to time $t$, the output equals that of $d$ at every time up to $t$, and the number of work
cells visited by time $t$ is at most $k(2r+1)$.
\end{lemma}
